\documentclass[UTF8,fontset=none,12pt]{ctexart}

\setCJKmainfont{Songti SC}
\setCJKsansfont{PingFang SC}
\setCJKmonofont{Songti SC}

\usepackage[a4paper,margin=2.35cm]{geometry}
\usepackage{amsmath,amssymb,amsthm,mathtools,bm}
\usepackage{graphicx}
\usepackage{float}
\usepackage{booktabs}
\usepackage{array}
\usepackage{enumitem}
\usepackage[hidelinks]{hyperref}

\title{数值线性代数\quad 第 24 讲：特征值问题}
\author{Lloyd N. Trefethen、David Bau III\quad 中文阅读版}
\date{}

\begin{document}
\maketitle
\pagestyle{plain}

特征值问题在科学计算中尤其有趣，因为寻找特征值的最佳算法虽然
功能强大，却往往非常不直观。本讲回顾特征值与特征向量的数学；
算法将在后面的讲次中讨论。

\section*{特征值与特征向量}

设 $A\in\mathbb{C}^{m\times m}$ 是方阵。若非零向量
$x\in\mathbb{C}^m$ 满足
\begin{equation}
Ax=\lambda x,
\tag{24.1}
\end{equation}
则称 $x$ 是 $A$ 的特征向量，$\lambda\in\mathbb{C}$ 是与之对应的
特征值。

这里的想法是，矩阵 $A$ 作用在 $\mathbb{C}^m$ 的某个子空间 $S$
上时，有时会表现得像标量乘法。发生这种情况时，这个特殊子空间
称为特征子空间，其中任意非零 $x\in S$ 都是特征向量。

矩阵 $A$ 的全部特征值组成 $A$ 的谱，记作
\[
\Lambda(A)\subset\mathbb{C}.
\]

特征值问题与前面几讲讨论的方程组问题有很不一样的性质。对方程组
而言，$A$ 的定义域可以是一个空间，值域可以是另一个空间。例如在
例 1.1 中，$A$ 把多项式系数的 $n$-向量映射为采样多项式值的
$m$-向量；询问这种 $A$ 的特征值没有意义。特征值问题只有在定义域
和值域相同的时候才有意义。这反映了应用中的一个事实：特征值通常用于
研究矩阵的迭代复合，或者研究 $A^k$ 这样的显式幂，或者研究
$e^{tA}$ 这样的函数形式。

概括地说，特征值和特征向量有两方面用途，一方面是算法上的，
另一方面是物理上的。从算法角度，特征值分析可以把某些耦合问题
化简为一组标量问题。从物理角度，特征值分析能够帮助我们理解
线性方程控制的演化系统的行为。最熟悉的例子包括共振（例如乐器
受到敲击、拨动或拉奏时的响应）和稳定性（例如流体受到微小扰动时
的行为）。在这些问题中，特征值对于分析大时间 $t$ 的行为尤其有用；
见习题 24.3。

\section*{特征值分解}

方阵 $A$ 的特征值分解（已经在 (5.1) 中提到）是如下因式分解：
\begin{equation}
A=X\Lambda X^{-1},
\tag{24.2}
\end{equation}
其中 $X$ 非奇异，$\Lambda$ 为对角矩阵。（如下文所述，这种分解
并不总是存在。）

这个定义可以改写为
\begin{equation}
AX=X\Lambda.
\tag{24.3}
\end{equation}
也就是
\[
A
\begin{bmatrix}x_1&x_2&\cdots&x_m\end{bmatrix}
=
\begin{bmatrix}x_1&x_2&\cdots&x_m\end{bmatrix}
\begin{bmatrix}
\lambda_1&&&\\
&\lambda_2&&\\
&&\ddots&\\
&&&\lambda_m
\end{bmatrix}.
\]
因此，如果 $x_j$ 是 $X$ 的第 $j$ 列、$\lambda_j$ 是 $\Lambda$ 的
第 $j$ 个对角条目，就有 $Ax_j=\lambda_jx_j$。换句话说，
$X$ 的第 $j$ 列是 $A$ 的特征向量，而 $\Lambda$ 的第 $j$ 个条目
是相应的特征值。

特征值分解把基变换为“特征向量坐标”。如果 $Ax=b$ 且
$A=X\Lambda X^{-1}$，则
\begin{equation}
X^{-1}b=\Lambda(X^{-1}x).
\tag{24.4}
\end{equation}
因此，为了计算 $Ax$，可以先在 $X$ 的列向量构成的基下展开 $x$，
再作用以 $\Lambda$，最后把结果理解为 $X$ 的列向量线性组合的系数。

\section*{几何重数}

对应于同一个特征值的全部特征向量，加上零向量后构成一个子空间，
称为特征子空间。若 $\lambda$ 是 $A$ 的特征值，记相应的特征子空间为
$E_\lambda$。特征子空间是 $A$ 的一个不变子空间，即
\[
AE_\lambda\subseteq E_\lambda.
\]

$E_\lambda$ 的维数可以理解为：对应于同一个特征值 $\lambda$ 的
线性无关特征向量的最大数量。这个数量称为 $\lambda$ 的几何重数。
几何重数也可以描述为
\[
\dim\mathcal{N}(A-\lambda I),
\]
因为这个零空间正是 $E_\lambda$。

\section*{特征多项式}

$A\in\mathbb{C}^{m\times m}$ 的特征多项式，记作 $p_A$，或简写为
$p$，是由下式定义的 $m$ 次多项式：
\begin{equation}
p_A(z)=\det(zI-A).
\tag{24.5}
\end{equation}
由于放置了负号，$p$ 是首一多项式，即 $m$ 次项的系数为 1。

\paragraph{定理 24.1.}
$\lambda$ 是 $A$ 的特征值，当且仅当 $p_A(\lambda)=0$。

\paragraph{证明.}
由特征值的定义，
\[
\begin{aligned}
\lambda\text{ 是 }A\text{ 的特征值}
&\Longleftrightarrow
\text{存在非零 }x\text{ 使 }\lambda x-Ax=0\\
&\Longleftrightarrow \lambda I-A\text{ 奇异}\\
&\Longleftrightarrow \det(\lambda I-A)=0.
\end{aligned}
\]
\hfill$\square$

定理 24.1 有一个重要推论：即使矩阵 $A$ 是实矩阵，它的某些特征值
也可能是复数。从物理上说，这与实动力系统既可能振荡，也可能增长
或衰减的现象有关。从算法上说，即使特征值问题的输入是实数，
输出也可能必须使用复数。

\section*{代数重数}

根据代数基本定理，可以把 $p_A$ 写成
\begin{equation}
p_A(z)=(z-\lambda_1)(z-\lambda_2)\cdots(z-\lambda_m)
\tag{24.6}
\end{equation}
的形式，其中 $\lambda_j\in\mathbb{C}$。由定理 24.1，每个
$\lambda_j$ 都是 $A$ 的特征值，$A$ 的所有特征值都会出现在这个列表中。
一般来说，一个特征值可能出现多次。特征值 $\lambda$ 的代数重数，
定义为它作为 $p_A$ 的根的重数。代数重数为 1 的特征值称为单特征值。

特征多项式给出了计算矩阵特征值数量的简单方法。

\paragraph{定理 24.2.}
若 $A\in\mathbb{C}^{m\times m}$，则 $A$ 有 $m$ 个特征值，
按代数重数计。特别地，如果 $p_A$ 的根都是单根，那么 $A$ 有
$m$ 个互不相同的特征值。

特别注意，每个矩阵至少有一个特征值。一个特征值的代数重数总是不小于
它的几何重数。为了证明这一点，需要先了解相似变换。

\section*{相似变换}

如果 $X\in\mathbb{C}^{m\times m}$ 非奇异，那么映射
\[
A\longmapsto X^{-1}AX
\]
称为 $A$ 的相似变换。如果存在非奇异的
$X\in\mathbb{C}^{m\times m}$ 使
\[
B=X^{-1}AX,
\]
就称矩阵 $A$ 和 $B$ 相似。正如特征值分解中的特殊情形所说明的，
相似变换就是一种换基操作。

许多性质会被相似矩阵 $A$ 和 $X^{-1}AX$ 共同保持。

\paragraph{定理 24.3.}
若 $X$ 非奇异，则 $A$ 和 $X^{-1}AX$ 具有相同的特征多项式、
特征值、代数重数和几何重数。

\paragraph{证明.}
特征多项式相同是直接计算的结果：
\[
\begin{aligned}
p_{X^{-1}AX}(z)
&=\det(zI-X^{-1}AX)\\
&=\det\bigl(X^{-1}(zI-A)X\bigr)\\
&=\det(X^{-1})\det(zI-A)\det(X)\\
&=\det(zI-A)=p_A(z).
\end{aligned}
\]
特征多项式相同，立即得到特征值和代数重数相同。最后，若
$E_\lambda$ 是 $A$ 的特征子空间，则可以验证
$X^{-1}E_\lambda$ 是 $X^{-1}AX$ 的特征子空间，反之亦然，
所以几何重数也相同。\hfill$\square$

现在可以把几何重数和代数重数联系起来。

\paragraph{定理 24.4.}
特征值的代数重数至少与它的几何重数一样大。

\paragraph{证明.}
设 $\lambda$ 对矩阵 $A$ 的几何重数为 $n$。令 $\widehat V$ 为
$m\times n$ 矩阵，其列构成特征子空间
\[
\{x:Ax=\lambda x\}
\]
的一组标准正交基。把 $\widehat V$ 扩充为一个方阵酉矩阵 $V$，
则
\[
B=V^*AV=
\begin{bmatrix}
\lambda I&C\\
0&D
\end{bmatrix},
\tag{24.7}
\]
其中 $I$ 是 $n\times n$ 单位矩阵，$C$ 是 $n\times(m-n)$ 矩阵，
$D$ 是 $(m-n)\times(m-n)$ 矩阵。由行列式的定义，
\[
\det(zI-B)=\det(zI-\lambda I)\det(zI-D)
=(z-\lambda)^n\det(zI-D).
\]
因此，$\lambda$ 作为 $B$ 的特征值，其代数重数至少为 $n$。
相似变换保持重数，所以对 $A$ 也成立。\hfill$\square$

\section*{缺陷特征值与缺陷矩阵}

虽然一般矩阵的代数重数和几何重数通常相等（也就是都为 1），
但这绝不是对所有矩阵都成立。

\paragraph{例 24.1.}
考虑两个矩阵
\[
A=
\begin{bmatrix}
2&0&0\\
0&2&0\\
0&0&2
\end{bmatrix},
\qquad
B=
\begin{bmatrix}
2&1&0\\
0&2&1\\
0&0&2
\end{bmatrix}.
\]
$A$ 和 $B$ 的特征多项式都是 $(z-2)^3$，因此二者都只有一个特征值
$\lambda=2$，代数重数为 3。对 $A$，可以选取三个线性无关的
特征向量，例如 $e_1,e_2,e_3$，所以几何重数也是 3。对 $B$，
却只能找到一个线性无关的特征向量（$e_1$ 的数乘），因此几何重数
只有 1。

代数重数大于几何重数的特征值称为缺陷特征值；含有一个或多个
缺陷特征值的矩阵称为缺陷矩阵。

任意对角矩阵都是非缺陷的。对这样的矩阵，一个特征值的代数重数
和几何重数都等于它在对角线上出现的次数。

\section*{可对角化性}

非缺陷矩阵恰好就是存在特征值分解 (24.2) 的矩阵类。

\paragraph{定理 24.5.}
$m\times m$ 矩阵 $A$ 非缺陷，当且仅当它有特征值分解
\[
A=X\Lambda X^{-1}.
\]

\paragraph{证明.}
若 $A=X\Lambda X^{-1}$，由定理 24.3，$\Lambda$ 与 $A$ 相似，
所以二者有相同的特征值和重数。由于 $\Lambda$ 是对角矩阵，
它是非缺陷的，从而 $A$ 也是非缺陷的。

反过来，非缺陷矩阵必须有 $m$ 个线性无关的特征向量：不同特征值
对应的特征向量必然线性无关，并且每个特征值能够贡献与其代数重数
一样多的线性无关特征向量。把这 $m$ 个特征向量作为矩阵 $X$ 的列，
$X$ 就是非奇异的，而且
\[
A=X\Lambda X^{-1}.
\]
\hfill$\square$

鉴于这个结论，非缺陷矩阵的另一个名称是可对角化矩阵。
一个可对角化矩阵 $A$ 是否在某种意义上“表现得像”它的对角等价矩阵
$\Lambda$？答案取决于所衡量的行为，以及特征向量矩阵 $X$ 的条件数。
如果 $X$ 高度病态，那么从 $A$ 转换为 $\Lambda$ 时可能丢失大量信息。
见第 34 讲的“谨慎提示：非正规性”。

\section*{行列式与迹}

$A\in\mathbb{C}^{m\times m}$ 的迹是它的对角线元素之和：
\[
\operatorname{tr}(A)=\sum_{j=1}^m a_{jj}.
\]
迹和行列式都与特征值有简单关系。

\paragraph{定理 24.6.}
矩阵 $A$ 的行列式 $\det(A)$ 和迹 $\operatorname{tr}(A)$，
分别等于 $A$ 的特征值之积和之和，特征值按代数重数计：
\begin{equation}
\det(A)=\prod_{j=1}^{m}\lambda_j,
\qquad
\operatorname{tr}(A)=\sum_{j=1}^{m}\lambda_j.
\tag{24.8}
\end{equation}

\paragraph{证明.}
由 (24.5) 和 (24.6)，
\[
\det(A)=(-1)^m\det(-A)=(-1)^mp_A(0)=\prod_{j=1}^{m}\lambda_j.
\]
这证明了第一个公式。对于第二个公式，由 (24.5)，$p_A$ 的
$z^{m-1}$ 项系数是 $A$ 的对角线元素之和的负值，即
$-\operatorname{tr}(A)$；另一方面，由 (24.6)，这个系数也是
$-\sum_{j=1}^{m}\lambda_j$。因此
$\operatorname{tr}(A)=\sum_{j=1}^{m}\lambda_j$。\hfill$\square$

\section*{酉对角化}

有时，$m\times m$ 矩阵 $A$ 不仅有 $m$ 个线性无关的特征向量，
而且这些特征向量可以选成相互正交的。在这种情况下，称 $A$ 是
酉可对角化的，即存在酉矩阵 $Q$ 使
\begin{equation}
A=Q\Lambda Q^*.
\tag{24.9}
\end{equation}
这个分解既是特征值分解，也是奇异值分解；区别只在于 $\Lambda$
的对角条目可能带有复数符号。

Hermitian 矩阵就是一类已经见过的酉可对角化矩阵。下面的结论是
定理 24.9 的推论。

\paragraph{定理 24.7.}
Hermitian 矩阵是酉可对角化的，而且它的特征值都是实数。

酉可对角化的矩阵不只有 Hermitian 矩阵。其他例子包括
反 Hermitian 矩阵、酉矩阵、循环矩阵，以及这些矩阵加上单位矩阵的
数倍。一般地，酉可对角化矩阵有一个优美的刻画。若
\[
A^*A=AA^*,
\]
就称 $A$ 为正规矩阵。下面的结果是一个著名定理。

\paragraph{定理 24.8.}
矩阵酉可对角化，当且仅当它是正规矩阵。

\section*{Schur 分解}

还有一种矩阵分解在数值分析中实际上最有用，因为包括缺陷矩阵在内的
所有矩阵都可以进行这种分解。矩阵 $A$ 的 Schur 分解是
\begin{equation}
A=QTQ^*,
\tag{24.10}
\end{equation}
其中 $Q$ 是酉矩阵，$T$ 是上三角矩阵。注意，由于 $A$ 和 $T$
相似，$A$ 的特征值必然出现在 $T$ 的对角线上。

\paragraph{定理 24.9.}
每个方阵 $A$ 都有 Schur 分解。

\paragraph{证明.}
对 $A$ 的维数 $m$ 使用归纳法。当 $m=1$ 时结论显然成立；
设 $m\ge2$。取 $A$ 的任意一个特征向量 $x$，对应特征值为 $\lambda$。
将 $x$ 归一化，并令它成为酉矩阵 $U$ 的第一列。正如 (24.7)，
容易验证
\[
U^*AU=
\begin{bmatrix}
\lambda&B\\
0&C
\end{bmatrix}.
\]
根据归纳假设，存在矩阵 $V$ 和上三角矩阵 $T$ 使
$C=VTV^*$。令
\[
Q=U
\begin{bmatrix}
1&0\\
0&V
\end{bmatrix}.
\]
这是一个酉矩阵，并且
\[
Q^*AQ=
\begin{bmatrix}
\lambda&BV\\
0&T
\end{bmatrix}.
\]
这就是所需的 Schur 分解。\hfill$\square$

\section*{揭示特征值的分解}

前面介绍了三种揭示特征值的分解，也就是把矩阵化为能够显式显示
特征值的形式：

\begin{itemize}[leftmargin=2.5em]
    \item 对角化 $A=X\Lambda X^{-1}$ 当且仅当 $A$ 非缺陷；
    \item 酉对角化 $A=Q\Lambda Q^*$ 当且仅当 $A$ 正规；
    \item 酉三角化（Schur 分解）$A=QTQ^*$ 对所有矩阵都存在。
\end{itemize}

为了计算特征值，我们将构造上述某一种分解。一般而言会使用 Schur
分解，因为它对所有矩阵都适用。此外，由于算法中使用了酉变换，
所得算法通常具有数值稳定性。如果 $A$ 正规，那么 Schur 形式就是
对角的；特别地，如果 $A$ 是 Hermitian 矩阵，我们可以在计算全过程
中利用这种对称性，把 $A$ 化为对角形式，而且运算量至多是一般矩阵
所需运算量的一半。

\section*{习题}

\begin{enumerate}[label=\textbf{24.\arabic*.},leftmargin=*]
    \item 对下列每个命题，证明它为真，或者给出反例说明它为假。
    除非另有说明，$A\in\mathbb{C}^{m\times m}$；本题用“ew”表示
    eigenvalue（特征值），这个缩写来自德语 Eigenwert，特征向量
    的对应缩写“ev”来自 Eigenvektor。
    \begin{enumerate}[label=(\alph*),leftmargin=2em]
        \item 如果 $\lambda$ 是 $A$ 的特征值且 $\mu\in\mathbb{C}$，
        那么 $\lambda-\mu$ 是 $A-\mu I$ 的特征值。
        \item 如果 $A$ 是实矩阵且 $\lambda$ 是 $A$ 的特征值，
        那么 $-\lambda$ 也是 $A$ 的特征值。
        \item 如果 $A$ 是实矩阵且 $\lambda$ 是 $A$ 的特征值，
        那么 $\overline{\lambda}$ 也是 $A$ 的特征值。
        \item 如果 $\lambda$ 是 $A$ 的特征值且 $A$ 非奇异，
        那么 $\lambda^{-1}$ 是 $A^{-1}$ 的特征值。
        \item 如果 $A$ 的所有特征值都是 0，那么 $A=0$。
        \item 如果 $A$ 是 Hermitian 矩阵且 $\lambda$ 是 $A$ 的特征值，
        那么 $|\lambda|$ 是 $A$ 的一个奇异值。
        \item 如果 $A$ 可对角化且它的所有特征值都相同，那么 $A$ 是对角矩阵。
    \end{enumerate}

    \item 下面是对任意 $m\times m$ 矩阵 $A$（不论对称与否）成立的
    Gerschgorin 定理：$A$ 的每个特征值至少落在复平面上的 $m$ 个
    圆盘中的一个内，其中第 $i$ 个圆盘的圆心为 $a_{ii}$，半径为
    $\sum_{j\ne i}|a_{ij}|$。此外，如果其中 $n$ 个圆盘构成一个
    与其余 $m-n$ 个圆盘不相交的连通区域，那么该区域内恰好有 $n$
    个 $A$ 的特征值。
    \begin{enumerate}[label=(\alph*),leftmargin=2em]
        \item 证明 Gerschgorin 定理的第一部分。（提示：设 $\lambda$
        是 $A$ 的任意特征值，$x$ 是对应特征向量，并把 $x$ 的最大条目
        归一化为 1。）
        \item 证明第二部分。（提示：把 $A$ 连续变形为对角矩阵，
        并利用矩阵特征值是其条目的连续函数这一事实。）
        \item 对下列矩阵的特征值，利用 Gerschgorin 定理给出估计：
        \[
        A=
        \begin{bmatrix}
        8&1&0\\
        1&4&\varepsilon\\
        0&\varepsilon&1
        \end{bmatrix},
        \qquad |\varepsilon|<1.
        \]
        \item 设法证明 $A$ 的最小特征值 $\lambda_3$ 满足更紧的界
        \[
        |\lambda_3-1|\le\varepsilon^2.
        \]
        （提示：考虑对角相似变换。）
    \end{enumerate}

    \item 设 $A$ 是一个 $10\times10$ 随机矩阵，其条目服从标准正态分布，
    再减去两倍单位矩阵。编写程序，在对数坐标下绘制
    $0\le t\le20$ 时的 $\|e^{tA}\|_2$，并将结果与直线
    $e^{t\alpha(A)}$ 比较，其中
    \[
    \alpha(A)=\max_j\operatorname{Re}(\lambda_j)
    \]
    是 $A$ 的谱 abscissa。对十个随机矩阵 $A$ 运行程序并评论结果。
    什么性质会使矩阵的 $\|e^{tA}\|_2$ 曲线在 $t\to\infty$ 时仍然振荡？

    \item 对任意 $A\in\mathbb{C}^{m\times m}$ 和任意范数，利用定理 24.9
    证明：
    \begin{enumerate}[label=(\alph*),leftmargin=2em]
        \item
        \[
        \lim_{n\to\infty}\|A^n\|=0
        \quad\Longleftrightarrow\quad
        \rho(A)<1,
        \]
        其中 $\rho$ 是谱半径（见习题 3.2）。
        \item
        \[
        \lim_{t\to\infty}\|e^{tA}\|=0
        \quad\Longleftrightarrow\quad
        \alpha(A)<0,
        \]
        其中 $\alpha$ 是谱 abscissa。
    \end{enumerate}
\end{enumerate}

\end{document}
